\documentclass{article}
\usepackage{preamble}

\begin{document}
\header{Filemon Mateus}{CS 6140}{Graphs \& Page Rank}{Homework \# 8}{\today}
\begin{enumerate}[leftmargin={*}, font={\bf}, label={\arabic*.}, ref={\arabic*}]
  \item \label{qst:1}
    \begin{enumerate}[label={(\Alph*)}, ref={\theenumi(\Alph*)}]
      \item
        Below is the implementation of a subroutine that for each node calculates
        the number of edges leading away from it (``out links'').
        \lstinputlisting
          [caption={\sf out-links.py}, label={lst:out-links}]
          {dev-dir/week8/out-links.py}

      \item
        Below is the implementation of a subroutine that fetches all the dead-end
        nodes in the graph but does not report the IDs of nodes that are dead-end
        nodes. There are in total {\bf 1235} in the graph.
        \lstinputlisting
          [caption={\sf dead-ends.py}, label={lst:dead-ends}]
          {dev-dir/week8/dead-ends.py}
          
      \newpage

      \item
        From the implementation below, the node with original ID {\bf 486980} has
        the highest {\sf PageRank} score.
        \lstinputlisting
          [caption={\sf pagerank-adj-mat.py}, label={lst:pagerank-adj-mat}]
          {dev-dir/week8/pagerank-adj-mat.py}

      \item
        From \autoref{lst:pagerank-adj-mat}, a total of {\bf 155} nodes point to the
        node with highest {\sf PageRank}. This is reasonable because it is the same
        output obtained from the {\sf PageRank} subroutine from
        \href{https://networkx.org/}{\url{NetworkX}} package.
    \end{enumerate}
    
  \newpage

  \item \label{qst:2}
    Below are implementations of some of the ways for speeding-up {\sf PageRank}.

    \begin{enumerate}[leftmargin={*}, label={\arabic*.}]
      \item
        Representing Transition Matrices.
        \lstinputlisting
          [caption={\sf pagerank-adj-lst.py}, label={lst:pagerank-adj-lst}]
          {dev-dir/week8/pagerank-adj-lst.py}
        Running \autoref{lst:pagerank-adj-lst} on the entire
        \href{https://snap.stanford.edu/data/web-Google.html}{\url{web-Google.txt}}
        dataset took {\bf 100} seconds. In contrast, \autoref{lst:pagerank-adj-mat}
        failed to run due to the sheer size of the adjacency matrix totaling
        $\SI{5}{\tera\byte}$ of disk space. This matrix representation is a performance
        improvement as the space required is roughly linear in the number of non-zero
        entries rather than quadratic in the size of the matrix.\footnote{
          Solution outline weakly inspired from:
          \href{
            http://infolab.stanford.edu/~ullman/mmds/ch5.pdf
          }{\url{http://infolab.stanford.edu/~ullman/mmds/ch5.pdf}}.
        }

      \item
        Distributed Map Reduce on $n$ Stripes. \par
        Below is \autoref{lst:pagerank-map-red}, a pseudo-implementation of
        {\sf PageRank} with MapReduce. It took {\bf 200} seconds to process the entire
        \href{https://snap.stanford.edu/data/web-Google.html}{\url{web-Google.txt}}
        dataset. This drop in performance is attributed to its sequential nature. In it,
        we observe a lack of concurrency where the {\tt map-reduce} function uses
        {\tt starmap} and {\tt reduce} in a serial manner, without leveraging parallelism
        or concurrency. In theory, since map-reduce is inherently parallelizable, running
        the code on multiple cores or nodes in a distributed cluster would significantly
        speed up computation and offer better performance.
        
        \newpage
        
        \lstinputlisting
          [caption={\sf pagerank-map-red.py}, label={lst:pagerank-map-red}]
          {dev-dir/week8/pagerank-map-red.py}
    \end{enumerate}
\end{enumerate}
\end{document}
